\Chapter{45}

\Verse{1}
{
    \zA{话说凤姐儿正抚恤平儿，忽见众姐妹进来，忙让了坐，平儿斟上茶来。 凤姐儿 笑道：“今儿来的这些人，倒像下帖子请了来的。” 探春先笑道：“我们有两件事：
        一件是我的，一件是四妹妹的，还夹着老太太的话。” 凤姐儿笑道：“有什么事这 么要紧？” 探春笑道：“我们起了个诗社，头一社就不齐全，众人脸软，所以就乱
        了例了。 我想必得你去做个‘监社御史’，铁面无私才好。}
}
{
    \eA{Lady Feng, we will now go on to explain, was engaged in comforting P'ing
        Erh, when upon unawares perceiving the young ladies enter the room, she
        hastened to make them sit down while P'ing Erh poured the tea. "So many of
        you come to-day," lady Feng smiled, "that it looks as if you'd been asked to
        come by invitation." T'an Ch'un was the first to speak.}
}
{
}

\Verse{2}
{
    \zA{再四妹妹为画园子，用 的东西这般那般不全，回了老太太，老太太说：‘只怕后头楼底下还有先剩下的， 找一找。 若有呢，拿出来； 若没有，叫人买去。’”
        凤姐儿笑道：“我又不会做什 么‘湿’咧‘干’的，叫我吃东西去倒会。” 探春笑道：“你不会做，也不用你做；
        你只监察着我们里头有偷安怠惰的，该怎么罚他就是了。” 凤姐儿笑道：“你们别 哄我，我早猜着了，那里是请我做‘监察御史’？ 分明叫了我去做个进钱的铜商罢
        咧。}
}
{
    \eA{"We have," she smilingly rejoined, "two objects in view, the one concerns
        me; the other cousin Quarta; but among these are, besides, certain things
        said by our venerable senior." "What's up?" inquired lady Feng with a laugh.
        "Is it so urgent?" "Some time ago," T'an Ch'un proceeded laughingly, "we
        started a rhyming club; but the first meeting was not quite a success. Every
        one of us proved so soft-hearted!}
}
{
}

\Verse{3}
{
    \zA{你们弄什么社，必是要轮流着做东道儿。 你们的钱不够花，想出这个法子来勾 了我去，好和我要钱。 可是这个主意不是？” 说的众人都笑道：“你猜着了！” 李
        纨笑道：“真真你是个水晶心肝玻璃人儿。” 凤姐笑道：“亏了你是个大嫂子呢! 姑娘们原是叫你带着念书，学规矩，学针线哪!这会子起诗社!能用几个钱，你就不
        管了？ 老太太、太太罢了，原是老封君。}
}
{
    \eA{The rules therefore were set at naught. So I can't help thinking that we
        must enlist your services as president of the society and superintendent;
        for what is needed to make the thing turn out well is firmness and no
        favour.}
}
{
}

\Verse{4}
{
    \zA{你一个月十两银子的月钱，比我们多两倍 子，老太太、太太还说你‘寡妇失业’的，可怜，不够用，又有个小子，足足的又 添了十两银子，和老太太、太太平等；
        又给你园子里的地，各人取租子； 年终分年 例，你又是上上分儿。 你娘儿们主子奴才共总没有十个人，吃的穿的仍旧是大官中 的。 通共算起来，也有四五百银子。
        这会子你就每年拿出一二百两来陪着他们玩玩 儿，有几年呢？ 他们明儿出了门子，难道你还赔不成？}
}
{
    \eA{The next matter is: cousin Quarta explained to our worthy ancestor that the
        requisites for painting the picture of the garden were short of one thing
        and another, and she said: 'that there must still be,' she fancied, 'in the
        lower story of the back loft some articles, remaining over from previous
        years, and that we should go and look for them.}
}
{
}

\Verse{5}
{
    \zA{这会子你怕花钱，挑唆他们 来闹我，我乐得去吃个河落海干，我还不知道呢！” 李纨笑道：“你们听听，我说了一句，他就说了两车无赖的话!真真泥腿光棍，
        专会打细算盘、分金掰两的。 你这个东西，亏了还托生在诗书仕宦人家做小姐，又 是这么出了嫁，还是这么着。 要生在贫寒小门小户人家，做了小子丫头，还不知怎
        么下作呢!天下人都叫你算计了去!昨儿还打平儿，亏你伸的出手来。 那黄汤难道灌 丧了狗肚子里去了？ 气的我只要替平儿打抱不平儿。}
}
{
    \eA{That if there be any, they should be taken out, but that in the event of
        their being none, some one should be commissioned to go and purchase a
        supply of them.'" "I'm not up to doing anything wet or dry, (play on word
        'shih,' verses)," lady Feng laughed, "and would you have me, pray, come and
        gorge?"}
}
{
}

\Verse{6}
{
    \zA{忖夺了半日，好容易‘狗长尾 巴尖儿’的好日子，又怕老太太心里不受用，因此没来。 究竟气还不平，你今儿倒 招我来了。
        给平儿拾鞋还不要呢!你们两个，很该换一个过儿才是。” 说的众人都 笑了。 凤姐忙笑道：“哦，我知道了，竟不是为诗为画来找我，竟是为平儿报仇来 了。
        我竟不知道平儿有你这么位仗腰子的人。 想来就像有鬼拉着我的手似的，从今 我也不敢打他了。 平姑娘，过来，我当着你大奶奶、姑娘们替你赔个不是，担待我
        ‘酒后无德’罢！”}
}
{
    \eA{"You may, it's possible, not be up to any of these things," T'an Ch'un
        replied, "but we don't expect you to do anything! All we want you for is to
        see whether there be among us any remiss or lazy, and to decide how they
        should be punished, that's all." "You shouldn't try and play your tricks
        upon me!" lady Feng smiled, "I can see through your little game! Is it that
        you wish me to act as president and superintendent?}
}
{
}

\Verse{7}
{
    \zA{说着众人都笑了。 李纨笑问平儿道：“如何？ 我说必要给你争 争气才罢。” 平儿笑道：“虽是奶奶们取笑儿，我可禁不起呢。” 李纨道：“什么
        禁的起禁不起，有我呢。 快拿钥匙叫你主子开门找东西去罢。” 凤姐儿笑道：“好嫂子!你且同他们去园子里去。 才要把这米帐合他们算一算，
        那边大太太又打发人来叫，又不知有什么话说，须得过去走一走。 还有你们年下添 补的衣裳，打点给人做去呢。”}
}
{
    \eA{No! it's as clear as day that your object is that I should play the part of
        that copper merchant, who put in contributions in hard cash. You have, at
        every meeting you hold, to each take turn and pay the piper; but, as your
        funds are not sufficient, you've invented this plan to come and inveigle me
        into your club, in order to wheedle money out of me! This must be your
        little conspiracy!"}
}
{
}

\Verse{8}
{
    \zA{李纨笑道：“这些事情我都不管，你只把我的事完 了，我好歇着去，省了这些姑娘们闹我。” 凤姐儿忙笑道：“好嫂子，赏我一点空 儿。
        你是最疼我的，怎么今儿为平儿就不疼我了？ 往常你还劝我说：‘事情虽多， 也该保全身子，检点着偷空儿歇歇。’ 你今儿倒反逼起我的命来了。 况且误了别人
        年下的衣裳无碍，他姐儿们的要误了，却是你的责任。 老太太岂不怪你不管闲事， 连一句现成的话也不说？ 我宁可自己落不是，也不敢累你呀。”}
}
{
    \eA{These words evoked general laughter. "You've guessed right!" they exclaimed.
        "In very truth," Li Wan smiled, "you're a creature with an intellect as
        transparent as crystal, and with wits as clear as glass!"}
}
{
}

\Verse{9}
{
    \zA{李纨笑道：“你们 听听，说的好不好？ 把他会说话的!我且问你：这诗社到底管不管？” 凤姐儿笑道： “这是什么话？
        我不入社花几个钱，我不成了大观园的反叛了么，还想在这里吃饭 不成？ 明日一早就到任，下马拜了印，先放下五十两银子给你们慢慢的做会社东道 儿。
        我又不会作诗作文的，只不过是个大俗人罢了。 ‘监察’也罢，不‘监察’也 罢，有了钱了，愁着你们还不撵出我来！” 说的众人又都笑起来。}
}
{
    \eA{"You've got the good fortune of being their elder sister-in-law," lady Feng
        smilingly remarked, "so the young ladies asked you to take them in hand, and
        teach them how to read, and make them learn good manners and needlework; and
        it's for you to guide and direct them in everything! But here they start a
        rhyming society, for which not much can be needed, and don't you concern
        yourself about them?}
}
{
}

\Verse{10}
{
    \zA{凤姐儿道：“过会子我开了楼房，所有这些东西，叫人搬出来你们瞧，要使得， 留着使； 要少什么，照你们的单子，我叫人赶着买去就是了。 画绢我就裁出来。 那
        图样没有在老太太那里，那边珍大爷收着呢。 说给你们，省了碰钉子去。 我去打发 人取了来，一并叫人连绢交给相公们矾去。 好不好呢？”
        李纨点头笑道：“这难为 你。 果然这么着还罢了。 那么着，咱们家去罢。 等着他不送了去，再来闹他。” 说 着便带了他姐妹们就走。}
}
{
    \eA{We'll leave our worthy ancestor and our Madame Wang aside; they are old
        people, but you receive each moon an allowance of ten taels, which is twice
        as much as what any one of us gets. More, our worthy ancestor and Madame
        Wang maintain that being a widow, and having lost your home, you haven't,
        poor thing, enough to live upon, and that you have a young child as well to
        bring up;}
}
{
}

\Verse{11}
{
    \zA{凤姐儿道：“这些事再没别人，都是宝玉生出来的。” 李 纨听了，忙回身笑道：“正为宝玉来，倒忘了他!头一社是他误了。 我们脸软，你 说该怎么罚他？”
        凤姐想了想，说道：“没别的法子，只叫他把你们各人屋子里的 地罚他扫一遍就完了。” 众人都笑道：“这话不差。”
        说着才要回去，只见一个小丫头扶着赖嬷嬷进来。 凤姐等忙站起来，笑道：“大 娘坐下。” 又都向他道喜。 赖嬷嬷向炕沿上坐了，笑道：“我也喜，主子们也喜。}
}
{
    \eA{so they added with extreme liberality another ten taels to your original
        share. Your allowance therefore is on a par with that of our dear senior.
        But they likewise gave you a piece of land in the garden, and you also come
        in for the lion's share of rents, collected from various quarters, and of
        the annual allowances, apportioned at the close of each year.}
}
{
}

\Verse{12}
{
    \zA{要不是主子们的恩典，我这喜打那里来呢？ 昨儿奶奶又打发彩哥赏东西，我孙子在 门上朝上磕了头了。” 李纨笑道：“多早晚上任去？” 赖嬷嬷叹道：“我那里管他
        们？ 由他们去罢。 前儿在家里给我磕头，我没好话。 我说：‘小子，别说你是官了，
        横行霸道的!你今年活了三十岁，虽然是人家的奴才，一落娘胎胞儿，主子的恩典， 放你出来，上托着主子的洪福，下托着你老子娘，也是公子哥儿似的读书写字，也
        是丫头、老婆、奶子捧凤凰似的。}
}
{
    \eA{Yet, you and your son don't muster, masters and servants, ten persons in
        all. What you eat and what your wear comes, just as ever, out of the general
        public fund, so that, computing everything together, you get as much as four
        to five hundred taels. Were you then to contribute each year a hundred or
        two hundred taels, to help them to have some fun, how many years could this
        outlay continue?}
}
{
}

\Verse{13}
{
    \zA{长了这么大，你那里知道那奴才两字是怎么写？ 只知道享福，也不知你爷爷和你老子受的那苦恼，熬了两三辈子，好容易挣出你这
        个东西，从小儿三灾八难，花的银子照样打出你这个银人儿来了。 到二十岁上，又 蒙主子的恩典，许你捐了前程在身上。 你看那正根正苗，忍饥挨饿的，要多少？ 你
        一个奴才秧子，仔细折了福!如今乐了十年，不知怎么弄神弄鬼，求了主子，又选 出来了。 县官虽小，事情却大，作那一处的官，就是那一方的父母。}
}
{
    \eA{They'll very soon be getting married, and, are they likely then to still
        expect you to make any contributions? So loth are you, however, at present
        to fork out any cash that you've egged them on to come and worry me! I'm
        quite prepared to spend away until we've drained our chest dry! Don't I know
        that the money isn't mine?" "Just you listen to her," Li Wan laughed.}
}
{
}

\Verse{14}
{
    \zA{你不安分守己， 尽忠报国，孝敬主子，只怕天也不容你。’” 李纨凤姐儿都笑道：“你也多虑。 我 们看他也就好。 先那几年，还进来了两次，这有好几年没来了。
        年下生日，只见他 的名字就罢了； 前儿给老太太、太太磕头来，在老太太那院里，见他又穿着新官的 服色，倒发的威武了，比先时也胖了。
        他这一得了官，正该你乐呢，反倒愁起这些 来!他不好，还有他的父母呢，你只受用你的就完了。 闲时坐个轿子进来，和老太
        太斗斗牌，说说话儿，谁好意思的委屈了你。}
}
{
    \eA{"I simply made one single remark, and out she came with two cartloads of
        nonsensical trash! You're as rough a diamond as a leg made of clay! All
        you're good for is to work the small abacus, to divide a catty and to
        fraction an ounce, so finicking are you! A nice thing you are, and yet,
        you've been lucky enough to come to life as the child of a family of learned
        and high officials.}
}
{
}

\Verse{15}
{
    \zA{家去一般也是楼房厦厅，谁不敬你？ 自然也是老封君似的了。” 平儿斟上茶来，赖嬷嬷忙站起来道：“姑娘不管叫那孩子倒来罢了，又生受你。”
        说着，一面吃茶，一面又道：“奶奶不知道，这小孩子们全要管的严。 饶这么严， 他们还偷空儿闹个乱子来，叫大人操心。 知道的，说小孩子们淘气； 不知道的，人
        家就说仗着财势欺人，连主子名声也不好。 恨的我没法儿，常把他老子叫了来，骂 一顿才好些。”}
}
{
    \eA{You've also made such a splendid match; and do you still behave in the way
        you do? Had you been a son or daughter born in some poverty-stricken, humble
        and low household, there's no saying what a mean thing you wouldn't have
        been! Every one in this world has been gulled by you; and yesterday you went
        so far as to strike P'ing Erh! But it wasn't the proper thing for you to
        stretch out your hand on her!}
}
{
}

\Verse{16}
{
    \zA{因又指宝玉道：“不怕你嫌我：如今老爷不过这么管你一管，老太 太就护在头里。 当日老爷小时，你爷爷那个打，谁没看见的!老爷小时，何曾像你
        这么天不怕地不怕的。 还有那边大老爷，虽然淘气，也没像你这扎窝子的样儿，也 是天天打。 还有东府里你珍大哥哥的爷爷，那才是火上浇油的性子，说声恼了，什
        么儿子，竟是审贼!如今我眼里看着，耳朵里听着，那珍大爷管儿子，倒也像当日 老祖宗的规矩，只是着三不着两的。}
}
{
    \eA{Was all that liquor, forsooth, poured down a cur's stomach? My monkey was
        up, and I meant to have taken upon myself to avenge P'ing Erh's grievance;
        but, after mature consideration, I thought to myself, 'her birthday is as
        slow to come round as a dog's tail grows to a point.' I also feared lest our
        venerable senior might be made to feel unhappy; so I did not come forward.
        Anyhow, my resentment isn't yet spent;}
}
{
}

\Verse{17}
{
    \zA{他自己也不管一管自己，这些兄弟侄儿怎么怨 的不怕他？ 你心里明白，喜欢我说； 不明白，嘴里不好意思，心里不知怎么骂我呢。”
        说着，只见赖大家的来了，接着周瑞家的张材家的都进来回事情。 凤姐儿笑道： “媳妇来接婆婆来了。” 赖大家的笑道：“不是接他老人家来的，倒是打听打听奶
        奶姑娘们赏脸不赏脸？” 赖嬷嬷听了，笑道：“可是我糊涂了!正经说的都没说， 且说些陈谷子烂芝麻的。 因为我们小子选出来了，众亲友要给他贺喜，少不得家里
        摆个酒。}
}
{
    \eA{and do you come to-day to try and irritate me? You aren't fit to even pick
        up shoes for P'ing Erh! You two should therefore change your respective
        places!" These taunts created merriment among the whole party. "Oh!" hastily
        exclaimed lady Feng, laughingly, "I know everything! You don't at all come
        to look me up on account of verses or paintings, but simply to take revenge
        on P'ing Erh's behalf!}
}
{
}

\Verse{18}
{
    \zA{我想摆一日酒，请这个不请那个也不是。 又想了一想，托主子的洪福，想 不到的这么荣耀光彩，就倾了家我也愿意的。 因此吩咐了他老子连摆三日酒：头一
        日在我们破花园子里摆几席酒，一台戏，请老太太、太太们、奶奶、姑娘们去散一 日闷，外头大厅上一台戏，几席酒，请老爷们、爷们，增增光； 第二日再请亲友；
        第三日再把我们两府里的伴儿请一请。 热闹三天，也是托着主子的洪福一场，光辉 光辉。” 李纨凤姐儿都笑道：“多早晚的日子？ 我们必去。}
}
{
    \eA{I never had any idea that P'ing Erh had such a backer as yourself to bolster
        her up! Had I known it, I wouldn't have ventured to strike her, even though
        a spirit had been tugging my arm! Miss P'ing come over and let me tender my
        apologies to you, in the presence of your senior lady and the young ladies.
        Do bear with me for having proved so utterly wanting in virtue, after I had
        had a few drinks!"}
}
{
}

\Verse{19}
{
    \zA{只怕老太太高兴要去也 定不得。” 赖大家的忙道：“择的日子是十四，只看我们奶奶的老脸罢了。” 凤姐 儿笑道：“别人我不知道，我是一定去的。
        先说下：我可没有贺礼，也不知道放赏， 吃了一走儿，可别笑话。” 赖大家的笑道：“奶奶说那里话？ 奶奶一喜欢，赏我们 三二万银子那就有了。”
        赖嬷嬷笑道：“我才去请老太太，老太太也说去，可算我 这脸还好。” 说毕叮咛了一回，方起身要走。}
}
{
    \eA{Every one felt amused by her insinuations. "What do you say?" Li Wan asked
        P'ing Erh smiling. "As for me, I think it my bounden duty to vindicate your
        wrongs, before we let the matter drop!" "Your remarks, ladies, may be spoken
        in jest," P'ing Erh smiled, "but I am not worthy of such a fuss!" "What
        about worthy and unworthy?" Li Wan observed. "I'm here for you!}
}
{
}

\Verse{20}
{
    \zA{因看见周瑞家的，便想起一事来，因 说道：“可是还有一句话问奶奶：这周嫂子的儿子，犯了什么不是，撵了他不用？”
        凤姐儿听了，笑道：“正是我要告诉你媳妇儿呢。 事情多，也忘了。 赖嫂子回去说 给你老头子，两府里不许收留他儿子，叫他各人去罢。” 赖大家的只得答应着。
        周瑞家的忙跪下央求。 赖嬷嬷忙道：“什么事？ 说给我评评。” 凤姐儿道：“前 儿我的生日，里头还没喝酒，他小子先醉了。}
}
{
    \eA{Quick, get the key, and let your mistress go and open the doors and hunt up
        the things!" "Dear sister-in-law," lady Feng said with a smile, "you'd
        better go along with them into the garden. I'm about to take the rice
        accounts in hand and square them up with them. Our senior lady, Madame
        Hsing, has also sent some one to call me; what she wants to tell me again, I
        can't make out; but I must need go over for a turn.}
}
{
}

\Verse{21}
{
    \zA{老娘那边送了礼来，他不在外头张罗， 倒坐着骂人，礼也不送进来。 两个女人进来了，他才带领小么儿们往里端。 小么儿
        们倒好好的，他拿的一盒子倒失了手，撒了一院子馒头。 人去了，我打发彩明去说 他，他倒骂了彩明一顿。 这样无法无天的忘八羔子，还不撵了做什么！”
        赖嬷嬷道： “我当什么事情，原来为这个。 奶奶听我说：他有不是，打他骂他，叫他改过就是 了； 撵出去断乎使不得。}
}
{
    \eA{There are, besides, all those extra clothes for you people to wear at the
        end of the year, and I must get them ready and give them to be made!" "These
        matters are none of my business!" Li Wan laughingly answered. "First settle
        my concerns so as to enable me to retire to rest, and escape the bother of
        having all these girls at me!"}
}
{
}

\Verse{22}
{
    \zA{他又比不得是咱们家的家生子儿，他现是太太的陪房，奶 奶只顾撵了他，太太的脸上不好看。 我说奶奶教导他几板子，以戒下次，仍旧留着 才是。
        不看他娘，也看太太。” 凤姐儿听了，便向赖大家的说道：“既这么着，明 儿叫了他来，打他四十棍，以后不许他喝酒。” 赖大家的答应了。 周瑞家的才磕头
        起来，又要给赖嬷嬷磕头，赖大家的拉着方罢。 然后他三人去了。 李纨等也就回园中来。 至晚，果然凤姐命人找了许多旧收的画具出来，送至园 中。
        宝钗等选了一回。}
}
{
    \eA{"Dear sister-in-law," vehemently smiled lady Feng, "be good enough to give
        me a little time! You've ever been the one to love me best, and how is it
        that you have, on P'ing Erh's account, ceased to care for me?}
}
{
}

\Verse{23}
{
    \zA{各色东西可用的只有一半，将那一半开了单子，给凤姐去照 样置买，不必细说。 一日外面矾了绢，起了稿子进来。 宝玉每日便在惜春那边帮忙，探春、李纨、
        迎春、宝钗等也都往那里来闲坐，一则观画，二则便于会面。 宝钗因见天气凉爽， 夜复渐长，遂至贾母房中商议，打点些针线来。 日间至贾母、王夫人处两次省候，
        不免又承色陪坐； 闲时园中姐妹处，也要不时闲话一回。 故日间不大得闲，每夜灯 下女工，必至三更方寝。}
}
{
    \eA{Time and again have you impressed on my mind that I should, despite my
        manifold duties, take good care of my health, and manage things in such a
        way as to find a little leisure for rest, and do you now contrariwise come
        to press the very life out of me? There's another thing besides. Should such
        clothes as will be required at the end of the year by any other persons be
        delayed, it won't matter;}
}
{
}

\Verse{24}
{
    \zA{黛玉每岁至春分、秋分后必犯旧疾，今秋又遇着贾母高兴， 多游玩了两次，未免过劳了神，近日又复嗽起来。 觉得比往常又重，所以总不出门， 只在自己房中将养。
        有时闷了，又盼个姐妹来说些闲话排遣； 及至宝钗等来望候他， 说不得三五句话，又厌烦了。 众人都体谅他病中，且素日形体娇弱，禁不得一些委
        屈，所以他接待不周，礼数疏忽，也都不责他。 这日宝钗来望他，因说起这病症来。}
}
{
    \eA{but, should those of the young ladies be behind time, let the responsibility
        rest upon your shoulders! And won't our old lady bear you a grudge, if you
        don't mind these small things? But as for me, I won't utter a single word
        against you, for, as I had rather bear the blame myself, I won't venture, to
        involve you!" "Listen to her!" Li Wan smiled. "Hasn't she got the gift of
        the gab? But let me ask you.}
}
{
}

\Verse{25}
{
    \zA{宝钗道：“这里走的几个大夫，虽都还好， 只是你吃他们的药，总不见效，不如再请一个高手的人来瞧一瞧，治好了岂不好？
        每年间闹一春一夏，又不老，又不小，成什么，也不是个常法儿。” 黛玉道：“不 中用。 我知道我的病是不能好的了。 且别说病，只论好的时候我是怎么个形景儿，
        就可知了。” 宝钗点头道：“可正是这话。 古人说，‘食谷者生’，你素日吃的竟 不能添养精神气血，也不是好事。” 黛玉叹道：“‘死生有命，富贵在天’，也不
        是人力可强求的。}
}
{
    \eA{Will you, after all, assume the control of this rhyming society or not?"
        "What's this nonsense you're talking?" lady Feng laughed. "Were I not to
        enter the society, and spend a little money, won't I be treated as a rebel
        in this garden of Broad Vista? And will I then still think of tarrying here
        to eat my head off?}
}
{
}

\Verse{26}
{
    \zA{今年比往年反觉又重了些似的。” 说话之间，已咳嗽了两三次。 宝钗道：“昨儿我看你那药方上，人参肉桂觉得太多了。 虽说益气补神，也不宜太 热。
        依我说：先以平肝养胃为要。 肝火一平，不能克土，胃气无病，饮食就可以养 人了。 每日早起，拿上等燕窝一两、冰糖五钱，用银铞子熬出粥来，要吃惯了，比
        药还强，最是滋阴补气的。” 黛玉叹道：“你素日待人，固然是极好的，然我最是个多心的人，只当你有心 藏奸。}
}
{
    \eA{So soon as the day dawns to-morrow, I'll arrive at my post, dismount from my
        horse, and, after kneeling before the seals, my first act will be to give
        fifty taels for you to quietly cover the expenses of your meetings. Yet
        after a few days, I shall neither indite any verses, nor write any
        compositions, as I am simply a rustic boor, nothing more! But it will be
        just the same whether I assume the direction or not;}
}
{
}

\Verse{27}
{
    \zA{从前日你说看杂书不好，又劝我那些好话，竟大感激你。 往日竟是我错了， 实在误到如今。 细细算来，我母亲去世的时候，又无姐妹兄弟，我长了今年十五岁，
        竟没一个人像你前日的话教导我。 怪不得云丫头说你好。 我往日见他赞你，我还不 受用； 昨儿我亲自经过，才知道了。 比如你说了那个，我再不轻放过你的；
        你竟不 介意，反劝我那些话：可知我竟自误了。 若不是前日看出来，今日这话，再不对你 说。}
}
{
    \eA{for after you pocket my money, there's no fear of your not driving me out of
        the place!" As these words dropped from her lips, one and all laughed again.
        "I'll now open the loft," proceeded lady Feng. "Should there be any of the
        articles you want, you can tell the servants to bring them out for you to
        look at them! If any will serve your purpose, keep them and use them.}
}
{
}

\Verse{28}
{
    \zA{你方才叫我吃燕窝粥的话，虽然燕窝易得，但只我因身子不好了，每年犯了这 病，也没什么要紧的去处； 请大夫，熬药，人参，肉桂，已经闹了个天翻地覆了，
        这会子我又兴出新文来，熬什么燕窝粥，老太太、太太、凤姐姐这三个人便没话， 那些底下老婆子丫头们，未免嫌我太多事了。 你看这里这些人，因见老太太多疼了
        宝玉和凤姐姐两个，他们尚虎视眈眈，背地里言三语四的，何况于我？ 况我又不是 正经主子，原是无依无靠投奔了来的，他们已经多嫌着我呢。}
}
{
    \eA{If any be short, I'll bid a servant go and purchase them according to your
        list. I'll go at once and cut the satin for the painting. As for the plan,
        it isn't with Madame Wang; it's still over there, at Mr. Chia Chen's. I tell
        you all this so that you should avoid going over to Madame Wang's and
        getting into trouble! But I'll go and depute some one to fetch it.}
}
{
}

\Verse{29}
{
    \zA{如今我还不知进退， 何苦叫他们咒我？” 宝钗道：“这么说，我也是和你一样。” 黛玉道：“你如何比我？ 你又有母亲， 又有哥哥。
        这里又有买卖地土，家里又仍旧有房有地。 你不过亲戚的情分，白住在 这里，一应大小事情又不沾他们一文半个，要走就走了。 我是一无所有，吃穿用度，
        一草一木，皆是和他们家的姑娘一样，那起小人岂有不多嫌的？” 宝钗笑道：“将 来也不过多费得一副嫁妆罢了，如今也愁不到那里。”}
}
{
    \eA{I'll direct also a servant to take the satin and give it to the gentlemen to
        size with alum; will this be all right?" Li Wan nodded her head by way of
        assent and smiled. "This will be putting you to much trouble and
        inconvenience," she said. "But we must really act as you suggest. Well in
        that case, go home all of you, and, if after a time, she doesn't send the
        thing round, you can come again and bully her."}
}
{
}

\Verse{30}
{
    \zA{黛玉听了不觉红了脸，笑道： “人家把你当个正经人，才把心里烦难告诉你听，你反拿我取笑儿！” 宝钗笑道： “虽是取笑儿，却也是真话。
        你放心，我在这里一日，我与你消遣一日。 你有什么 委屈烦难，只管告诉我，我能解的，自然替你解。 我虽有个哥哥，你也是知道的； 只有个母亲，比你略强些。
        咱们也算同病相怜。 你也是个明白人，何必作‘司马牛 之叹’？ 你才说的也是，多一事不如省一事。 我明日家去和妈妈说了，只怕燕窝我
        们家里还有，与你送几两。}
}
{
    \eA{So saying, she there and then led off the young ladies, and was making her
        way out, when lady Feng exclaimed: "It's Pao-yü and he alone, who has given
        rise to all this fuss." Li Wan overheard her remark and hastily turned
        herself round. "We did, in fact, come over," she smiled, "on account of
        Pao-yü, and we forgot, instead all about him! The first meeting was deferred
        through him;}
}
{
}

\Verse{31}
{
    \zA{每日叫丫头们就熬了，又便宜，又不惊师动众的。” 黛 玉忙笑道：“东西是小，难得你多情如此。” 宝钗道：“这有什么放在嘴里的!只
        愁我人人跟前失于应候罢了。 这会子只怕你烦了，我且去了。” 黛玉道：“晚上再 来和我说句话儿。” 宝钗答应着便去了，不在话下。
        这里黛玉喝了两口稀粥，仍歪在床上。 不想日未落时，天就变了，淅淅沥沥下 起雨来。 秋霖脉脉，阴晴不定，那天渐渐的黄昏时候了，且阴的沉黑，兼着那雨滴
        竹梢，更觉凄凉。}
}
{
    \eA{but we are too soft-hearted, so tell us what penalty to inflict on him!"
        Lady Feng gave herself to reflection. "There's only one thing to do," she
        then remarked. "Just punish him by making him sweep the floor of each of
        your rooms. This will do!" "Your verdict is faultless!" they laughed with
        one accord.}
}
{
}

\Verse{32}
{
    \zA{知宝钗不能来了，便在灯下随便拿了一本书，却是《乐府杂稿》， 有《秋闺怨》、《别离怨》等词。 黛玉不觉心有所感，不禁发于章句，遂成《代别
        离》一首，拟《春江花月夜》之格，乃名其词为《秋窗风雨夕》。 词曰： 秋花惨淡秋草黄，耿耿秋灯秋夜长。 已觉秋窗秋不尽，那堪风雨助秋凉! 助秋风雨来何速？
        惊破秋窗秋梦续。 抱得秋情不忍眠，自向秋屏挑泪烛。 泪烛摇摇？ 短檠，牵愁照恨动离情。 谁家秋院无风入？ 何处秋窗无雨声？}
}
{
    \eA{While they conversed they were on the point of starting on their way back,
        when they caught sight of a young maid walk in, supporting nurse Lai. Lady
        Feng and her companions immediately rose to their feet, their faces beaming
        with smiles. "Venerable mother!" they said, "do take a seat!" They then in a
        body presented their congratulations to her.}
}
{
}

\Verse{33}
{
    \zA{罗衾不奈秋风力，残漏声催秋雨急。 连宵脉脉复飕飕，灯前似伴离人泣。 寒烟小院转萧条，疏竹虚窗时滴沥。 不知风雨几时休，已教泪洒窗纱湿。
        吟罢搁笔，方欲安寝，丫鬟报说：“宝二爷来了。” 一语未尽，只见宝玉头上 戴着大箬笠，身上披着蓑衣。 黛玉不觉笑道：“那里来的这么个渔翁？” 宝玉忙问：
        “今儿好？ 吃了药了没有？ 今儿一日吃了多少饭？” 一面说，一面摘了笠，脱了蓑。}
}
{
    \eA{Nurse Lai seated herself on the edge of the stovecouch and returned their
        smiles. "I'm to be congratulated," she rejoined, "but you, mistresses, are
        to be congratulated as well; for had it had not been for the bountiful grace
        displaced by you, mistresses, whence would this joy of mine have come?}
}
{
}

\Verse{34}
{
    \zA{一手举起灯来，一手遮着灯儿，向黛玉脸上照了一照，觑着瞧了一瞧，笑道：“今 儿气色好了些。” 黛玉看他脱了蓑衣，里面只穿半旧红绫短袄，系着绿汗巾子，膝
        上露出绿绸撒花裤子，底下是掐金满绣的绵纱袜子，？ 着蝴蝶落花鞋。 黛玉问道： “上头怕雨，底下这鞋袜子是不怕的？ 也倒干净些呀。”
        宝玉笑道：“我这一套是 全的。 一双棠木屐，才穿了来，脱在廊檐下了。” 黛玉又看那蓑衣斗笠不是寻常市 卖的，十分细致轻巧，因说道：“是什么草编的？}
}
{
    \eA{Your ladyship sent Ts'ai Ko again yesterday to bring me presents, but my
        grandson \_kotowed\_ at the door, with his face turned towards the upper
        quarters." "When is he going to his post?" Li Wan inquired, with a smile.
        Nurse Lai heaved a sigh. "How can I interfere with them?" she answered.
        "Why, I let them have their own way and start when they like!}
}
{
}

\Verse{35}
{
    \zA{怪道穿上不像那刺猬似的。” 宝 玉道：“这三样都是北静王送的。 他闲常下雨时，在家里也是这样。 你喜欢这个， 我也弄一套来送你。
        别的都罢了，惟有这斗笠有趣：上头这顶儿是活的，冬天下雪 戴上帽子，就把竹信子抽了去，拿下顶子来，只剩了这个圈子，下雪时男女都带得。
        我送你一顶，冬天下雪戴。” 黛玉笑道：“我不要他。 戴上那个，成了画儿上画的 和戏上扮的那渔婆儿了。”}
}
{
    \eA{The other day, they were at my house, and they prostrated themselves before
        me; but I could find no complimentary remark to make to him, so, 'Sir!' I
        said, 'putting aside that you're an official, you've lived in a reckless and
        dissolute way, for now thirty years.}
}
{
}

\Verse{36}
{
    \zA{及说了出来，方想起来这话恰与方才说宝玉的话相连了， 后悔不迭，羞的脸飞红，伏在桌上，嗽个不住。
        宝玉却不留心，因见案上有诗，遂拿起来看了一遍，又不觉叫好。 黛玉听了， 忙起来夺在手内，灯上烧了。 宝玉笑道：“我已记熟了。” 黛玉道：“我要歇了，
        你请去罢，明日再来。” 宝玉听了，回手向怀内掏出一个核桃大的金表来，瞧了一 瞧，那针已指到戌末亥初之间，忙又揣了，说道：“原该歇了，又搅的你劳了半日
        神。”}
}
{
    \eA{You should, it's true, have been people's bond-servant, but from the moment
        you came out of your mother's womb, your master graciously accorded you your
        liberty. Thanks, above, to the boundless blessings showered upon you by your
        lord, and, below, to the favour of your father and mother, you're like a
        noble scion and a gentleman, able to read and to write;}
}
{
}

\Verse{37}
{
    \zA{说着，披蓑戴笠出去了，又翻身进来，问道：“你想什么吃？ 你告诉我，我 明儿一早回老太太，岂不比老婆子们说的明白？” 黛玉笑道：“等我夜里想着了，
        明日一早告诉你。 你听雨越发紧了，快去罢。 可有人跟没有？” 两个婆子答应：“有， 在外面拿着伞点着灯笼呢。” 黛玉笑道：“这个天点灯笼？”
        宝玉道：“不相干， 是羊角的，不怕雨。” 黛玉听说，回手向书架上把个玻璃绣球灯拿下来，命点一枝 小蜡儿来，递与宝玉道：“这个又比那个亮，正是雨里点的。”}
}
{
    \eA{and you have been carried about by maids, old matrons, and nurses, just as
        if you had been a very phoenix! But now that you've grown up and reached
        this age, do you have the faintest notion of what the two words
        'bond-servant' imply? All you think of is to enjoy your benefits.}
}
{
}

\Verse{38}
{
    \zA{宝玉道：“我也有 这么一个，怕他们失脚滑倒了打破了，所以没点来。” 黛玉道：“跌了灯值钱呢， 是跌了人值钱？ 你又穿不惯木屐子。
        那灯笼叫他们前头点着，这个又轻巧又亮，原 是雨里自己拿着的。 你自己手里拿着这个，岂不好？ 明儿再送来。 就失了手也有限
        的，怎么忽然又变出这‘剖腹藏珠’的脾气来！” 宝玉听了，随过来接了。 前头两 个婆子打着伞，拿着羊角灯，后头还有两个小丫鬟打着伞。}
}
{
    \eA{But what hardships your grandfather and father had to bear, in slaving away
        for two or three generations, before they succeeded, after ever so many ups
        and downs, in raising up a thing like you, you don't at all know! From your
        very infancy, you ever ailed from this, or sickened for that, so that the
        money that was expended on your behalf, would suffice to fuse into a
        lifelike silver image of you!}
}
{
}

\Verse{39}
{
    \zA{宝玉便将这个灯递给一 个小丫头捧着，宝玉扶着他的肩，一径去了。 就有蘅芜院两个婆子，也打着伞提着灯，送了一大包燕窝来，还有一包子洁粉 梅片雪花洋糖。
        说：“这比买的强。 我们姑娘说：‘姑娘先吃着，完了再送来。’” 黛玉回说：“费心。” 命他：“外头坐了吃茶。” 婆子笑道：“不喝茶了，我们还 有事呢。”
        黛玉笑道：“我也知道你们忙。 如今天又凉，夜又长，越发该会个夜局， 赌两场了。” 一个婆子笑道：“不瞒姑娘说，今年我沾了光了。}
}
{
    \eA{At the age of twenty, you again received the bounty of your master in the
        shape of a promise to purchase official status for you. But just mark, how
        many inmates of the principal branch and main offspring have to endure
        privation, and suffer the pangs of hunger! So beware you, who are the
        offshoot of a bond-servant, lest you snap your happiness!}
}
{
}

\Verse{40}
{
    \zA{横竖每夜有几个上 夜的人，误了更又不好，不如会个夜局，又坐了更，又解了闷。 今儿又是我的头家， 如今园门关了，就该上场儿了。”
        黛玉听了，笑道：“难为你们。 误了你们的发财， 冒雨送来。” 命人：“给他们几百钱打些酒吃，避避雨气。” 那两个婆子笑道：“又 破费姑娘赏酒吃。”
        说着磕了头，出外面接了钱，打伞去了。 紫鹃收起燕窝，然后移灯下帘，伏侍黛玉睡下。 黛玉自在枕上感念宝钗，一时 又羡他有母有兄；
        一回又想宝玉素昔和睦，终有嫌疑。}
}
{
    \eA{After enjoying so many good things for a decade, by the help of what
        spirits, and the agency of what devils have you, I wonder, managed to so
        successfully entreat your master as to induce him to bring you to the fore
        again and select you for office? Magistrates may be minor officials, but
        their functions are none the less onerous.}
}
{
}

\Verse{41}
{
    \zA{又听见窗外竹梢蕉叶之上， 雨声淅沥，清寒透幕，不觉又滴下泪来。 直到四更方渐渐的睡熟了。 暂且无话。 要知端底，且看下回分解。 (本章完)}
}
{
    \eA{In whatever district they obtain a post, they become the father and mother
        of that particular locality. If you therefore don't mind your business, and
        look after your duties in such a way as to acquit yourself of your loyal
        obligations, to prove your gratitude to the state and to show obedience and
        reverence to your lord, heaven, I fear, will not even bear with you!'" Li
        Wan and lady Feng laughed.}
}
{
}

\Verse{42}
{
    \zA{}
}
{
    \eA{"You're too full of misgivings!" they observed. "From what we can see of
        him, he's all right! Some years back, he paid us a visit or two; but it's
        many years now that he hasn't put his foot here. At the close of each year,
        and on birthdays, we've simply seen his name brought in, that's all. The
        other day, that he came to knock his head before our venerable senior and
        Madame Wang, we caught sight of him in her courtyard yonder; and, got up in
        the uniform of his new office, he looked so dignified, and stouter too than
        before. Now that he has got this post, you should be quite happy; instead of
        that you worry and fret about this and that!}
}
{
}

\Verse{43}
{
    \zA{}
}
{
    \eA{If he does get bad, why, he has his father and mother yet to take care of
        him, so all you need do is to be cheerful and content! When you've got time
        to spare, do get into a chair and come in and have a game of cards and a
        chat with our worthy senior; and who ever will have the face to hurt your
        feelings? Why, were you go to your home, you'd also have there houses and
        halls, and who is there who would not hold you in high respect? You're
        certainly, what one would call, a venerable old dame!" P'ing Erh poured a
        cup of tea and brought it to her. Nurse Lai speedily stood up. "You could
        have asked any girl to do this for me; it wouldn't have mattered! But here
        I'm troubling you again!"}
}
{
}

\Verse{44}
{
    \zA{}
}
{
    \eA{Apologising, she resumed, sipping her tea the while: "My lady you're not
        aware that young girls of this age must be in everything kept strictly in
        hand. In the event of any license, they're sure to find time to kick up
        trouble, and annoy their elders. Those, who know (how well they are
        supervised), will then say that children are always up to mischief. But
        those, who don't, will maintain that they take advantage of their wealthy
        position to despise people; to the detriment as well of their mistresses'
        reputation. How I regret that there's nothing that I can do with him. Time
        after time, have I had to send for his father; and he has been the better,
        after a scolding from him."}
}
{
}

\Verse{45}
{
    \zA{}
}
{
    \eA{Pointing at Pao-yü, "I don't mind whether you feel angry with me for what
        I'm going to say," she proceeded, "but if your father were to attempt now to
        exercise ever so little control over you, your venerable grandmother is sure
        to try and screen you. Yet, when in days gone by your worthy father was
        young, he used to be beaten by your grandfather. Who hasn't seen him do it?
        But did your father, in his youth resemble you, who have neither fear for
        God or man? There was also our senior master, on the other side, Mr. Chia
        She. He was, I admit, wild; but never such a crossgrained fellow as
        yourself; and yet he too had his daily dose of the whip. There was besides
        the father of your elder cousin Chen, of the eastern mansion.}
}
{
}

\Verse{46}
{
    \zA{}
}
{
    \eA{He had a disposition that flared up like a fire over which oil is poured. If
        anything was said, and he flew into a rage, why, talk about a son, it was
        really as if he tortured a robber. From all I can now see and hear, Mr. Chen
        keeps his son in check just as much as was the custom in old days among his
        ancestors; the only thing is that he abides by it in some respects, but not
        in others. Besides, he doesn't exercise the least restraint over his own
        self, so is it to be wondered at if all his cousins and nieces don't respect
        him? If you've got any sense about you, you'll only be too glad that I speak
        to you in this wise; but if you haven't, you mayn't be very well able to say
        anything openly to me, but you'll inwardly abuse me, who knows to what
        extent!"}
}
{
}

\Verse{47}
{
    \zA{}
}
{
    \eA{As she reproved him, they saw Lai Ta's wife arrive. In close succession came
        Chou Jui's wife along with Chang Ts'ai's wife to report various matters. "A
        wife," laughed lady Feng, "has come to fetch her mother-in-law!" "I haven't
        come to fetch our old dame," Lai Ta's wife smilingly rejoined, "but to
        inquire whether you, my lady and the young ladies, will confer upon us the
        honour of your company?" When nurse Lai caught this remark, she smiled.
        "I've really grown quite idiotic!" "What," she exclaimed, "was right and
        proper for me to say, I didn't say, but I went on talking instead a lot of
        rot and rubbish! As our relatives and friends are presenting their
        congratulations to our grandson for having been selected to fill up that
        office of his, we find ourselves under the necessity of giving a banquet at
        home.}
}
{
}

\Verse{48}
{
    \zA{}
}
{
    \eA{But I was thinking that it wouldn't do, if we kept a feast going the whole
        day, and we invited this one, and not that one. Reflecting also that it was
        thanks to our master's vast bounty that we've come in for this unforeseen
        glory and splendour, I felt quite agreeable to do anything, even though it
        may entail the collapse of our household. I therefore advised his father to
        give banquets on three consecutive days. That he should, on the first, put
        up several tables, and a stage in our mean garden, and invite your venerable
        dowager lady, the senior ladies, junior ladies, and young ladies to come and
        have some distraction during the day, and that he should have several tables
        laid on the stage in the main pavilion outside, and request the senior and
        junior gentlemen to confer upon us the lustre of their presence.}
}
{
}

\Verse{49}
{
    \zA{}
}
{
    \eA{That for the second day, we should ask our relatives and friends; and that
        for the third, we should invite our companions from the two mansions. In
        this way, we'll have three days' excitement, and, by the boundless favour of
        our master, we'll have the benefit of enjoying the honour of your society."
        "When is it to be?" Li Wan and lady Feng inquired, smilingly. "As far as we
        are concerned, we'll feel it our duty to come. And we hope that our worthy
        senior may feel in the humour to go. But there's no saying for certain!"
        "The day chosen is the fourteenth," Lai Ta's wife eagerly replied. "Just
        come for the sake of our old mother-in-law!" "I can't tell about the
        others," lady Feng explained with a laugh, "but as for me I shall positively
        come.}
}
{
}

\Verse{50}
{
    \zA{}
}
{
    \eA{I must however tell you beforehand that I've no congratulatory presents to
        give you. Nor do I know anything about tips to players or others. As soon as
        I shall have done eating, I shall bolt, so don't laugh at me."
        "Fiddlesticks!" Lai Ta's wife laughed. "Were your ladyship disposed, you
        could well afford to give us twenty and thirty thousand taels." "I'm off now
        to invite our venerable mistress," nurse Lai smilingly remarked. "And if her
        ladyship also agrees to come, I shall deem it a greater honour than ever
        conferred upon me." Having said this, she went on to issue some injunctions;
        after which, she got up to go, when the sight of Chou Jui's wife reminded
        her of something. "Of course!" she consequently observed. "I've got one more
        question to ask you, my lady.}
}
{
}

\Verse{51}
{
    \zA{}
}
{
    \eA{What did sister-in-law Chou's son do to incur blame, that he was packed off,
        and his services dispensed with?" "I was just about to tell your
        daughter-in-law," lady Feng answered smilingly, after listening to her
        question, "but with so many things to preoccupy me, it slipped from my
        memory! When you get home, sister-in-law Lai, explain to that old husband of
        yours that we won't have his, (Chou Jui's), son kept in either of the
        mansions; and that he can tell him to go about his own business!" Lai Ta's
        wife had no option but to express her acquiescence. Chou Jui's wife however
        speedily fell on her knees and gave way to urgent entreaties. "What is it
        all about?" nurse Lai shouted.}
}
{
}

\Verse{52}
{
    \zA{}
}
{
    \eA{"Tell me and let me determine the right and wrong of the question." "The
        other day," lady Feng observed, "that my birthday was celebrated, that young
        fellow of his got drunk, before the wine ever went round; and when the old
        dame, over there, sent presents, he didn't go outside to give a helping
        hand, but squatted down, instead, and upbraided people. Even the presents he
        wouldn't carry inside. And it was only after the two girls had come indoors
        that he eventually got the servant-lads and brought them in. Those lads were
        however careful enough in what they did, but as for him, he let the box, he
        held, slip from his hands, and bestrewed the whole courtyard with cakes.
        When every one had left, I deputed Ts'ai Ming to go and talk to him; but he
        then turned round and gave Ts'ai Ming a regular scolding.}
}
{
}

\Verse{53}
{
    \zA{}
}
{
    \eA{So what's the use of not bundling off a disorderly rascal like him, who
        neither shows any regard for discipline or heaven?" "I was wondering what it
        could be!" nurse Lai ventured. "Was it really about this? My lady, listen to
        me! If he has done anything wrong, thrash him and scold him, until you make
        him mend his ways, and finish with it! But to drive him out of the place,
        will never, by any manner of means, do. He isn't, besides, to be treated
        like a child born in our household. He is at present employed as Madame
        Wang's attendant, so if you carry out your purpose of expelling him, her
        ladyship's face will be put to the blush. My idea is that you should, my
        lady, give him a lesson by letting him have several whacks with a cane so as
        to induce him to abstain from wine in the future.}
}
{
}

\Verse{54}
{
    \zA{}
}
{
    \eA{If you then retain him in your service as hitherto he'll be all right! If
        you don't do it for his mother's sake; do it at least for that of Madame
        Wang!" After lending an ear to her arguments, lady Feng addressed herself to
        Lai Ta's wife. "Well, in that case," she said, "call him over to-morrow and
        give him forty blows; and don't let him after this touch any more wine!" Lai
        Ta's wife promised to execute her directions. Chou Jui's wife then kotowed
        and rose to her feet. But she also persisted upon prostrating herself before
        nurse Lai; and only desisted when Lai Ta's wife pulled her up. But presently
        the trio took their departure, and Li Wan and her companions sped back into
        the garden.}
}
{
}

\Verse{55}
{
    \zA{}
}
{
    \eA{When evening came, lady Feng actually bade the servants go and look (into
        the loft), and when they discovered a lot of painting materials, which had
        been put away long ago, they brought them into the garden. Pao-ch'ai and her
        friends then selected such as they deemed suitable. But as they only had as
        yet half the necessaries they required, they drew out a list of the other
        half and sent it to lady Feng, who, needless for us to particularise, had
        the different articles purchased, according to the specimens supplied. By a
        certain day, the silk had been sized outside, a rough sketch drawn, and both
        returned into the garden.}
}
{
}

\Verse{56}
{
    \zA{}
}
{
    \eA{Pao-yü therefore was day after day to be found over at Hsi Ch'un's, doing
        his best to help her in her hard work. But T'an Ch'un, Li Wan, Ying Ch'un,
        Pao-ch'ai and the other girls likewise congregated in her quarters, and sat
        with her when they were at leisure, as they could, in the first place, watch
        the progress of the painting, and as secondly they were able to conveniently
        see something of each other. When Pao-ch'ai perceived how cool and pleasant
        the weather was getting, and how the nights were beginning again to
        gradually draw out, she came and found her mother, and consulted with her,
        until they got some needlework ready.}
}
{
}

\Verse{57}
{
    \zA{}
}
{
    \eA{Of a day, she would cross over to the quarters of dowager lady Chia and
        Madame Wang, and twice pay her salutations, but, she could not help as well
        amusing them and sitting with them to keep them company. When free, she
        would come and see her cousins in the garden, and have, at odd times, a chat
        with them, so having, during daylight no leisure to speak of, she was wont,
        of a night, to ply her needle by lamplight, and only retire to sleep after
        the third watch had come and gone.}
}
{
}

\Verse{58}
{
    \zA{}
}
{
    \eA{As for Tai-yü, she had, as a matter of course, a relapse of her complaint
        regularly every year, soon after the spring equinox and autumn solstice. But
        she had, during the last autumn, also found her grandmother Chia in such
        buoyant spirits, that she had walked a little too much on two distinct
        occasions, and naturally fatigued herself more than was good for her.
        Recently, too, she had begun to cough and to feel heavier than she had done
        at ordinary times, so she never by any chance put her foot out of doors, but
        remained at home and looked after her health. When at times, dullness crept
        over her, she longed for her cousins to come and chat with her and dispel
        her despondent feelings.}
}
{
}

\Verse{59}
{
    \zA{}
}
{
    \eA{But whenever Pao-ch'ai or any of her cousins paid her a visit, she barely
        uttered half a dozen words before she felt quite averse to any society. Yet
        one and all made every allowance for her illness. And as she had ever been
        in poor health and not strong enough to resist any annoyance, they did not
        find the least fault with her, despite even any lack of propriety she showed
        in playing the hostess with them, or any remissness on her part in observing
        the prescribed rules of etiquette. Pao-ch'ai came, on this occasion to call
        on her. The conversation started on the symptoms of her ailment.}
}
{
}

\Verse{60}
{
    \zA{}
}
{
    \eA{"The various doctors, who visit this place," Pao-ch'ai consequently
        remarked, "may, it's true, be all very able practitioners; but you take
        their medicines and don't reap the least benefit! Wouldn't it be as well
        therefore to ask some other person of note to come and see you? And could he
        succeed in getting you all right, wouldn't it be nice? Here you year by year
        ail away throughout the whole length of spring and summer; but you're
        neither so old nor so young, so what will be the end of it? Besides, it
        can't go on for ever." "It's no use," Tai-yü rejoined. "I know well enough
        that there's no cure for this complaint of mine! Not to speak of when I'm
        unwell, why even when I'm not, my state is such that one can see very well
        that there's no hope!"}
}
{
}

\Verse{61}
{
    \zA{}
}
{
    \eA{Pao-ch'ai shook her head. "Quite so!" she ventured. "An old writer says:
        'Those who eat, live.' But what you've all along eaten hasn't been enough to
        strengthen your energies and physique. This isn't a good thing!" Tai-yü
        heaved a sigh. "Whether I'm to live or die is all destiny!" she said.
        "Riches and honours are in the hands of heaven; and human strength cannot
        suffice to forcibly get even them! But my complaint this year seems to be
        far worse than in past years, instead of any better." While deploring her
        lot, she coughed two or three times. "It struck me," Pao-ch'ai said, "that
        in that prescription of yours I saw yesterday there was far too much ginseng
        and cinnamon. They are splendid tonics, of course, but too many heating
        things are not good.}
}
{
}

\Verse{62}
{
    \zA{}
}
{
    \eA{I think that the first urgent thing to do is to ease the liver and give tone
        to the stomach. When once the fire in the liver is reduced, it will not be
        able to overcome the stomach; and, when once the digestive organs are free
        of ailment, drink and food will be able to give nutriment to the human
        frame. As soon as you get out of bed, every morning, take one ounce of
        birds' nests, of superior quality, and five mace of sugar candy and prepare
        congee with them in a silver kettle. When once you get into the way of
        taking this decoction, you'll find it far more efficacious than medicines;
        for it possesses the highest virtue for invigorating the vagina and bracing
        up the physique."}
}
{
}

\Verse{63}
{
    \zA{}
}
{
    \eA{"You've certainly always treated people with extreme consideration," sighed
        Tai-yü, "but such a supremely suspicious person am I that I imagined that
        you inwardly concealed some evil design! Yet ever since the day on which you
        represented to me how unwholesome it was to read obscene books, and you gave
        me all that good advice, I've felt most grateful to you! I've hitherto, in
        fact, been mistaken in my opinion; and the truth of the matter is that I
        remained under this misconception up to the very present. But you must
        carefully consider that when my mother died, I hadn't even any sisters or
        brothers; and that up to this my fifteenth year there has never been a
        single person to admonish me as you did the other day. Little wonder is it
        if that girl Yün speaks well of you!}
}
{
}

\Verse{64}
{
    \zA{}
}
{
    \eA{Whenever, in former days, I heard her heap praise upon you, I felt uneasy in
        my mind, but, after my experiences of yesterday, I see how right she was.
        When you, for instance, began to tell me all those things, I didn't forgive
        you at the time, but, without worrying yourself in the least about it you
        went on, contrariwise, to tender me the advice you did. This makes it
        evident that I have laboured under a mistaken idea! Had I not made this
        discovery the other day, I wouldn't be speaking like this to your very face
        to-day. You told me a few minutes back to take bird's nest congee; but
        birds' nests are, I admit, easily procured;}
}
{
}

\Verse{65}
{
    \zA{}
}
{
    \eA{yet all on account of my sickly constitution and of the relapses I have
        every year of this complaint of mine, which amounts to nothing, doctors have
        had to be sent for, medicines, with ginseng and cinnamon, have had to be
        concocted, and I've given already such trouble as to turn heaven and earth
        topsy-turvey; so were I now to start again a new fad, by having some birds'
        nests congee or other prepared, our worthy senior, Madame Wang, and lady
        Feng, will, all three of them, have no objection to raise; but that posse of
        matrons and maids below will unavoidably despise me for my excessive
        fussiness! Just notice how every one in here ogles wildly like tigers their
        prey;}
}
{
}

\Verse{66}
{
    \zA{}
}
{
    \eA{and stealthily says one thing and another, simply because they see how fond
        our worthy ancestor is of both Pao-yü and lady Feng, and how much more won't
        they do these things with me? What's more, I'm not a pucker mistress. I've
        really come here as a mere refugee, for I had no one to sustain me and no
        one to depend upon. They already bear me considerable dislike; so much so,
        that I'm still quite at a loss whether I should stay or go; and why should I
        make them heap execrations upon me?" "Well, in that case," Pao-ch'ai
        observed, "I'm too in the same plight as yourself!" "How can you compare
        yourself with me?" Tai-yü exclaimed. "You have a mother; and a brother as
        well! You've also got some business and land in here, and, at home, you can
        call houses' and fields your own.}
}
{
}

\Verse{67}
{
    \zA{}
}
{
    \eA{It's only therefore the ties of relationship, which make you stay here at
        all. Neither are you in anything whether large or small, in their debt for
        one single cash or even half a one; and when you want to go, you're at
        liberty to go. But I, have nothing whatever that I can call my own. Yet, in
        what I eat, wear, and use, I am, in every trifle, entirely on the same
        footing as the young ladies in their household, so how ever can that mean
        lot not despise me out and out?" "The only extra expense they'll have to go
        to by and bye," Pao-ch'ai laughed, "will be to get one more trousseau,
        that's all. And for the present, it's too soon yet to worry yourself about
        that!"}
}
{
}

\Verse{68}
{
    \zA{}
}
{
    \eA{At this insinuation, Tai-yü unconsciously blushed scarlet. "One treats you,"
        she smiled, "as a decent sort of person, and confides in you the woes of
        one's heart, and, instead of sympathising with me, you make me the means of
        raising a laugh!" "Albeit I raise a laugh at your expense," Pao-ch'ai
        rejoined, a smile curling her lips, "what I say is none the less true! But
        compose your mind! I'll try every day that I'm here to cheer you up; so come
        to me with every grievance or trouble, for I shall, needless to say, dispel
        those that are within my power. Notwithstanding that I have a brother, you
        yourself know well enough what he's like! All I have is a mother, so I'm
        just a trifle better off than you! We can therefore well look upon ourselves
        as being in the same boat, and sympathise with each other.}
}
{
}

\Verse{69}
{
    \zA{}
}
{
    \eA{You have, besides, plenty of wits about you, so why need you give way to
        groans, as did Ssu Ma-niu? What you said just now is quite right; but, you
        should worry and fret about as little and not as much as you can. On my
        return home, to-morrow, I'll tell my mother; and, as I think there must be
        still some birds' nests in our house, we'll send you several ounces of them.
        You can then tell the servant-maids to prepare some for you at whatever time
        you want every day; and you'll thus be suiting your own convenience and be
        giving no trouble or annoyance to any one." "The things are, of themselves,
        of little account," eagerly responded Tai-yü laughingly. "What's difficult
        to find is one with as much feeling as yourself." "What's there in this
        worth speaking about?" Pao-ch'ai said.}
}
{
}

\Verse{70}
{
    \zA{}
}
{
    \eA{"What grieves me is that I fail to be as nice as I should be with those I
        come across. But, I presume, you feel quite done up now, so I'll be off!"
        "Come in the evening again," Tai-yü pressed her, "and have a chat with me."
        While assuring her that she would come, Pao-ch'ai walked out, so let us
        leave her alone for the present. Tai-yü, meanwhile, drank a few sips of thin
        congee, and then once more lay herself down on her bed. But before the sun
        set, the weather unexpectedly changed, and a fine drizzling rain set in. So
        gently come the autumn showers that dull and fine are subject to uncertain
        alternations. The shades of twilight gradually fell on this occasion. The
        heavens too got so overcast as to look deep black.}
}
{
}

\Verse{71}
{
    \zA{}
}
{
    \eA{Besides the effect of this change on her mind, the patter of the rain on the
        bamboo tops intensified her despondency, and, concluding that Pao-ch'ai
        would be deterred from coming, she took up, in the lamp light, the first
        book within her reach, which turned out to be the 'Treasury of Miscellaneous
        Lyrics.' Finding among these 'the Pinings of a maiden in autumn,' 'the
        Anguish of Separation,' and other similar poems, Tai-yü felt unawares much
        affected; and, unable to restrain herself from giving vent to her feelings
        in writing, she, there and then, improvised the following stanza, in the
        same strain as the one on separation; complying with the rules observed in
        the 'Spring River-Flower' and 'Moonlight Night.'}
}
{
}

\Verse{72}
{
    \zA{}
}
{
    \eA{These verses, she then entitled 'the Poem on the Autumn evening, when wind
        and rain raged outside the window.' Their burden was:  In autumn, flowers
        decay; herbage, when autumn comes, doth yellow  turn. On long autumnal
        nights, the autumn lanterns with bright radiance  burn. As from my window
        autumn scenes I scan, autumn endless doth seem. This mood how can I bear,
        when wind and rain despondency enhance? How sudden break forth wind and
        rain, and help to make the  autumntide! Fright snaps my autumn dreams, those
        dreams which under my lattice I  dreamt. A sad autumnal gloom enclasps my
        heart, and drives all sleep away! In person I approach the autumn screen to
        snuff the weeping wick. The tearful candles with a flickering flame consume
        on their short  stands.}
}
{
}

\Verse{73}
{
    \zA{}
}
{
    \eA{They stir up grief, dazzle my eyes, and a sense of parting arouse. In what
        family's courts do not the blasts of autumn winds intrude? And where in
        autumn does not rain patter against the window-frames? The silken quilt
        cannot ward off the nipping force of autumn winds. The drip of the half
        drained water-clock impels the autumn rains. A lull for few nights reigned,
        but the wind has again risen in  strength. By the lantern I weep, as if I
        sat with some one who must go. The small courtyard, full of bleak mist, is
        now become quite  desolate. With quick drip drops the rain on the distant
        bamboos and vacant  sills. What time, I wonder, will the wind and rain their
        howl and patter  cease? The tears already I have shed have soakèd through
        the window gauze.}
}
{
}

\Verse{74}
{
    \zA{}
}
{
    \eA{After scanning her verses, she flung the pen aside, and was just on the
        point of retiring to rest, when a waiting-maid announced that 'master
        Secundus, Mr. Pao-yü, had come.' Barely was the announcement out of her
        lips, than Pao-yü appeared on the scene with a large bamboo hat on his head,
        and a wrapper thrown over his shoulders. Of a sudden, a smile betrayed
        itself on Tai-yü's lips. "Where does this fisherman come from?" she
        exclaimed. "Are you better to-day?" Pao-yü inquired with alacrity. "Have you
        had any medicines? How much rice have you had to eat to-day?" While plying
        her with questions, he took off the hat and divested himself of the wrapper;
        and, promptly raising the lamp with one hand, he screened it with the other
        and threw its rays upon Tai-yü's face.}
}
{
}

\Verse{75}
{
    \zA{}
}
{
    \eA{Then straining his eyes, he scrutinised her for a while. "You look better
        to-day," he smiled. As soon as he threw off his wrapper, Tai-yü noticed that
        he was clad in a short red silk jacket, the worse for wear; that he was
        girded with a green sash, and that, about his knees, his nether garments
        were visible, made of green thin silk, brocaded with flowers. Below these,
        he wore embroidered gauze socks, worked all over with twisted gold thread,
        and a pair of shoes ornamented with butterflies and clusters of fallen
        flowers. "Above, you fight shy of the rain," Tai-yü remarked, "but aren't
        these shoes and socks below afraid of rain? Yet they're quite clean!" "This
        suit is complete!" Pao-yü smiled. "I've got a pair of crab-wood clogs, I put
        on to come over;}
}
{
}

\Verse{76}
{
    \zA{}
}
{
    \eA{but I took them off under the eaves of the verandah." Tai-yü's attention was
        then attracted by the extreme fineness and lightness of the texture of his
        wrapper and hat, which were unlike those sold in the market places. "With
        what grass are they plaited?" she consequently asked. "It would be strange
        if you didn't, with this sort of things on, look like a very hedgehog!"
        "These three articles are a gift from the Prince of Pei Ching," Pao-yü
        answered. "Ordinarily, when it rains, he too wears this kind of outfit at
        home. But if it has taken your fancy, I'll have a suit made for you. There's
        nothing peculiar about the other things, but this hat is funny! The crown at
        the top is movable;}
}
{
}

\Verse{77}
{
    \zA{}
}
{
    \eA{so if you want to wear a hat, during snowy weather in wintertime, you pull
        off the bamboo pegs, and remove the crown, and there you only have the
        circular brim. This is worn, when it snows, by men and women alike. I'll
        give you one therefore to wear in the wintry snowy months." "I don't want
        it!" laughed Tai-yü. "Were I to wear this sort of thing, I'd look like one
        of those fisherwomen, one sees depicted in pictures or represented on the
        stage!"}
}
{
}

\Verse{78}
{
    \zA{}
}
{
    \eA{Upon reaching this point, she remembered that there was some connection
        between her present remarks and the comparison she had some time back made
        with regard to Pao-yü, and, before she had time to indulge in regrets, a
        sense of shame so intense overpowered her that the colour rushed to her
        face, and, leaning her head on the table, she coughed and coughed till she
        could not stop. Pao-yü, however, did not detect her embarrassment; but
        catching sight of some verses lying on the table, he eagerly snatched them
        up and conned them from beginning to end. "Splendid!" he could not help
        crying. But the moment Tai-yü heard his exclamation, she speedily jumped to
        her feet, and clutched the verses and burnt them over the lamp. "I've
        already committed them sufficiently to memory!" Pao-yü laughed.}
}
{
}

\Verse{79}
{
    \zA{}
}
{
    \eA{"I want to have a little rest," Tai-yü said, "so please get away; come back
        again to-morrow." At these words, Pao-yü drew back his hand, and producing
        from his breast a gold watch about the size of a walnut, he looked at the
        time. The hand pointed between eight and nine p.m.; so hastily putting it
        away, "You should certainly retire to rest!" he replied. "My visit has upset
        you. I've quite tired you out this long while." With these apologies, he
        threw the wrapper over him, put on the rain-hat and quitted the room. But
        turning round, he retraced his steps inside. "Is there anything you fancy to
        eat?" he asked. "If there be, tell me, and I'll let our venerable ancestor
        know of it to-morrow as soon as it's day. Won't I explain things clearer
        than any of the old matrons could?"}
}
{
}

\Verse{80}
{
    \zA{}
}
{
    \eA{"Let me," rejoined Tai-yü smiling, "think in the night. I'll let you know
        early to-morrow. But harken, it's raining harder than it did; so be off at
        once! Have you got any attendants, or no?" "Yes!" interposed the two
        matrons. "There are servants to wait on him. They're outside holding his
        umbrella and lighting the lanterns." "Are they lighting lanterns with this
        weather?" laughed Tai-yü. "It won't hurt them!" Pao-yü answered. "They're
        made of sheep's horn, so they don't mind the rain." Hearing this, Tai-yü put
        back her hand, and, taking down an ornamented glass lantern in the shape of
        a ball from the book case, she asked the servants to light a small candle
        and bring it to her; after which, she handed the lantern to Pao-yü. "This,"
        she said, "gives out more light than the others; and is just the thing for
        rainy weather."}
}
{
}

\Verse{81}
{
    \zA{}
}
{
    \eA{"I've also got one like it." Pao-yü replied. "But fearing lest they might
        slip, fall down and break it, I did not have it lighted and brought round."
        "What's of more account," Tai-yü inquired, "harm to a lantern or to a human
        being? You're not besides accustomed to wearing clogs, so tell them to walk
        ahead with those lanterns. This one is as light and handy as it is
        light-giving; and is really adapted for rainy weather, so wouldn't it be
        well if you carried it yourself? You can send it over to me to-morrow! But,
        were it even to slip from your hand, it wouldn't matter much. How is it that
        you've also suddenly developed this money-grabbing sort of temperament? It's
        as bad as if you ripped your intestines to secrete pearls in."}
}
{
}

\Verse{82}
{
    \zA{}
}
{
    \eA{After these words, Pao-yü approached her and took the lantern from her.
        Ahead then advanced two matrons, with umbrellas and sheep horn lanterns, and
        behind followed a couple of waiting-maids also with umbrellas. Pao-yü handed
        the glass lantern to a young maid to carry, and, supporting himself on her
        shoulder, he straightway wended his steps on his way back. But presently
        arrived an old servant from the Heng Wu court, provided as well with an
        umbrella and a lantern, to bring over a large bundle of birds' nests, and a
        packet of foreign sugar, pure as powder, and white as petals of plum-blossom
        and flakes of snow. "These," she said, "are much better than what you can
        buy. Our young lady sends you word, miss, to first go on with these.}
}
{
}

\Verse{83}
{
    \zA{}
}
{
    \eA{When you've done with them, she'll let you have some more." "Many thanks for
        the trouble you've taken!" Tai-yü returned for answer; and then asked her to
        go and sit outside and have a cup of tea. "I won't have any tea," the old
        servant smiled. "I've got something else to attend to." "I'm well aware that
        you've all got plenty in hand," Tai-yü resumed with a smiling countenance.
        "But the weather being cool now and the nights long, it's more expedient
        than ever to establish two things: a night club and a gambling place." "I
        won't disguise the fact from you, miss," the old servant laughingly
        observed, "that I've managed this year to win plenty of money. Several
        servants have, under any circumstances, to do night duty;}
}
{
}

\Verse{84}
{
    \zA{}
}
{
    \eA{and, as any neglect in keeping watch wouldn't be the right thing, isn't it
        as well to have a night club, as one can sit on the look-out and dispel
        dullness as well? But it's again my turn to play the croupier to-day, so I
        must be getting along to the place, as the garden gate, will, by this time,
        be nearly closing!" This rejoinder evoked a laugh from Tai-yü. "I've given
        you all this bother," she remarked, "and made you lose your chances of
        getting money, just to bring these things in the rain." And calling a
        servant she bade her present her with several hundreds of cash to buy some
        wine with, to drive the damp away. "I've uselessly put you again, miss, to
        the expense of giving me a tip for wine," the old servant smiled.}
}
{
}

\Verse{85}
{
    \zA{}
}
{
    \eA{But saying this she knocked her forehead before her; and issuing outside,
        she received the money, after which, she opened her umbrella, and trudged
        back. Tzu Chüan meanwhile put the birds' nests away; and removing afterwards
        the lamps, she lowered the portières and waited upon Tai-yü until she lay
        herself down to sleep. While she reclined all alone on her pillow, Tai-yü
        thought gratefully of Pao-ch'ai. At one moment, she envied her for having a
        mother and a brother; and at another, she mused that with the friendliness
        Pao-yü had ever shown her they were bound to be the victims of suspicion.
        But the pitter-patter of the rain, dripping on the bamboo tops and banana
        leaves, fell on her ear; and, as a fresh coolness penetrated the curtain,
        tears once more unconsciously trickled down her cheeks.}
}
{
}

\Verse{86}
{
    \zA{}
}
{
    \eA{In this frame of mind, she continued straight up to the fourth watch, when
        she at last gradually dropped into a sound sleep. For the time, however,
        there is nothing that we can add. So should you, reader, desire to know any
        subsequent details, peruse what is written in the next chapter.}
}
{
}
